\subsection{The genus-two norm obstruction and its exact local-matching scope}

\begin{result}{E.70}{Proved}{The complete projected Selmer group and the actual order-thirty-two quotient}
Put
\[
 C:y^2=x^5-54x,\quad J=\operatorname{Jac}(C),\quad
 K=\Q(\theta),\quad\theta^4=54,\quad
 L=\Q(r),\quad r^2=6,\quad\theta^2=3r.
\]
The cyclic-eight cover
\[
 W:\ z^8=-t(1-t)^3/24\longrightarrow U
 \tag{E.70.1}
\]
has proper compactification $C$. Its proper mod-two homology produces
an actual arithmetic quotient of $E_U$ with geometric group
$J[2]\rtimes C_8$, of order $128$ and nilpotency class four.
Writing $v$ for the cyclic deck generator and $\delta=v-1$, the
characteristic quotient by $\delta^2J[2]$ has geometric group
\[
 \overline Q=\overline M\rtimes C_8,\qquad
 \overline M=J[2]/\delta^2J[2]
       \simeq\operatorname{Res}_{L/\Q}\mu_2,
 \qquad |\overline Q|=32,\quad\operatorname{cl}(\overline Q)=2.
 \tag{E.70.2}
\]
Its continuous cohomology map is the quadratic norm
$K^\times/K^{\times2}\to L^\times/L^{\times2}$.
If $\mathcal S\subset H^1(G_\Q,\overline M)$ consists of the classes
whose localization at every place belongs to the projected Kummer
image of $J(\Q_w)/2J(\Q_w)$, then
\[
 \mathcal S=\langle[-(2+r)],[3+r]\rangle\simeq(\Z/2)^2.
 \tag{E.70.3}
\]
These four classes are the images of rational divisor classes.
In combined eighth-power class $[-24]$ for $t(1-t)^3$, the
everywhere locally point-or-cusp-matching section locus of this
order-$32$ arithmetic quotient with ordinary centre in
$828+1129\Z_{1129}$ is empty. This independently confirms the
finite quotient, projected Selmer calculation, and exclusion in
SECTIONS\_OVER\_Q A.209, without assuming its earlier full
two-Selmer computation A.186.
\end{result}
\begin{proof}
The mutually inverse function-field maps are
\[
 t=\frac{x^4}{x^4-54},\qquad z=\frac{3x}{y},\qquad
 x=-\frac{6z^2}{1-t},\qquad y=-\frac{18z}{1-t}.
 \tag{E.70.4}
\]
Substitution verifies both curve equations and both compositions.
The ramification indices over $0,1,\infty$ are respectively
$8,8,2$; there are four points over infinity. Thus the six boundary
points give $2g(C)-2=-16+7+7+4=2$.

Let $H\subset\Pi$ be the geometric index-eight subgroup of this
cover. Compactification and the Abel--Jacobi map based at the rational
point $\infty\in C(\Q)$ give a continuous surjection
\[
 H\longrightarrow\pi_1(C_{\overline\Q})^{\rm ab}
                    \longrightarrow J[2].
 \tag{E.70.5}
\]
Its kernel is preserved by deck transformations and arithmetic
automorphisms. It is consequently normal in $E_U$, so every quotient
below is a quotient of the actual arithmetic extension.

For the geometric calculation choose peripheral generators $x_0,y_0$,
put $v=x_0$, $u=y_0x_0^{-3}$, and use the Schreier generators
$u_j=v^juv^{-j}$, $0\le j<8$, and $z_0=v^8$ for $H$.
The zero-boundary relation kills $z_0$. The one-boundary word
$y_0^8$ has abelianization $\sum_j u_j+3z_0$.
Since $x_0y_0=u_1v^4$, its square gives $u_1+u_5+z_0$.
Its four conjugates give exactly $u_{j+4}=u_j$ in mod-two
proper homology; the sum relation is then redundant. Hence
\[
 \Pi/\ker(\text{E.70.5})
   =\bigl(\F_2[\delta]/\delta^4\bigr)\rtimes C_8,
 \qquad v\cdot m=(1+\delta)m.
 \tag{E.70.6}
\]
The successive lower-central orders are $128,8,4,2,1$.
After dividing by $\delta^2J[2]$ they are $32,2,1$.
In particular the latter kernel contains $\Gamma_3\Pi$, giving
the continuous arithmetic quotient map from the integral class-two
extension. The four-dimensional coefficient module is the actual
$J[2]$ from (E.70.5), not just an abstract group of the same order.

To identify its Galois action, order the four nonzero branch roots
as $\theta_j=i^j\theta$. Put
$T=[(0,0)-\infty]$, $T_j=[(\theta_j,0)-\infty]$ and
$e_j=T+T_j$. The standard branch description and Weil pairing
\cite[\S4, pp.~5--6]{RootGrossHanoi} give
$e_2(e_j,T_k)=(-1)^{\delta_{jk}}$; thus the $e_j$ form a basis.
The deck action cycles them, and
\[
 (a_0,a_1,a_2,a_3)\longmapsto(a_0+a_2,a_1+a_3)
 \tag{E.70.7}
\]
is the quotient by $\delta^2$. The two coordinates are indexed by
the two values of $\theta^2$. Every arithmetic permutation of the
four roots acts on these pairs through its action on $L$.
This proves the Galois-equivariant identification in (E.70.2).
In the branch Kummer description the four square-root coordinates
are $x-\theta_j$. Multiplying opposite coordinates is exactly
the norm $N_{K/L}$, which proves the cohomology map by Kummer
theory and continuous Shapiro. The fifth branch coordinate is
recovered as the norm from $K$ to $\Q$; it introduces no scalar
quotient. In particular, for $D=[(3,9)-\infty]$,
\[
 [P-\infty]\longmapsto[x(P)^2-3r],\quad
 D\longmapsto[9-3r],\quad T\longmapsto[-3r].
 \tag{E.70.8}
\]
The last evaluation uses the four nonzero branch functions, which
are units at $(0,0)$, so the vanishing fifth coordinate causes no
ambiguity.

We compute all of $\mathcal S$. The curve has good reduction
away from two and three. Properness extends every local Jacobian
point there to its abelian scheme, and multiplication by two is
finite etale. Its Kummer class, and then its projected class, is
unramified. Every candidate in $L^\times/L^{\times2}$ therefore
has even valuation outside the primes over two and three.
The integer ring is $\Z[r]$, of discriminant $24$.
Minkowski's bound is $\sqrt{24}/2<3$
\cite[Theorem~4.3, p.~70]{CohMilneANT}. The only possible
nontrivial ideal representative has norm two, and that prime is
generated by $\pi_2=2+r$, whose norm is $-2$. Thus the class
group is zero. Set $\pi_3=3+r$ and $\epsilon=5+2r$; direct
calculation gives
\[
 N(\pi_3)=3,\qquad \pi_2^2=2\epsilon,\qquad
 \pi_3^2=3\epsilon,\qquad N(\epsilon)=1.
 \tag{E.70.9}
\]
The unit group is $\langle-1,\epsilon\rangle$. Indeed norm
minus one is impossible modulo three. A positive norm-one unit
$a+br$ strictly between $1$ and $\epsilon$ would have
$0<b<2$ by $2b\sqrt6=u-u^{-1}$; $b=1$ would force
$a^2=7$. Multiplying by powers of $\epsilon$, and then changing
sign or taking the inverse, treats every unit.
Removing even prime valuations by global squares now gives
exactly the sixteen possible classes
\[
             (-1)^a\epsilon^b\pi_2^c\pi_3^d,
             \qquad a,b,c,d\in\F_2.
 \tag{E.70.10}
\]
Their independence follows from the two prime valuations and
the proved unit group.

At the real embedding $r=-\sqrt6$ the norm $K/L$ is the
complex-to-real norm, and is positive. The signs in (E.70.10)
force $a=c$. At three, $X^4-54$ is irreducible over $\Q_3$:
each root has valuation $3/4$, so a monic factor of degree $j$
would have constant coefficient of valuation $3j/4$, forcing
$4\mid j$. The branch permutation module consequently has
two invariant elements, so $|J(\Q_3)[2]|=2$.
An open formal pro-three subgroup has invertible multiplication
by two. On the remaining finite group its kernel and cokernel
have equal order; hence $|J(\Q_3)/2J(\Q_3)|=2$.

In the ramified quadratic field $L_3$, normalize $v(r)=1$,
so $v(3)=2$. The projected class $[-3r]$ has odd valuation
and is nonzero. It is therefore the whole nonzero local image.
The exact identities
\[
 9-3r=9/\pi_3,\qquad (-3r)/(9-3r)=-\pi_2
 \tag{E.70.11}
\]
show that it is the class of $\pi_3$: the unit $-\pi_2$
has residue one and is a square by simple-root Hensel lifting.
The unit $\epsilon$ has nonsquare residue two; together with
the odd-valuation $\pi_3$ these are a basis for the local
squareclasses. Thus $b=0$. Precisely the four classes
$(-\pi_2)^a\pi_3^d$ remain. Conversely (E.70.8) and
(E.70.11) realize all four by the rational divisor classes
$D,T$. They satisfy every local condition, including the
dyadic one. This proves the all-place assertion (E.70.3).

At $p=1129$ the two roots of $r^2=6$ are $838,291$.
Both global generators are squares in both factors:
\[
\begin{array}{c|cc}
 &r=838&r=291\\\hline
 -\pi_2&289=17^2&836=426^2\\
 \pi_3&841=29^2&294=221^2
\end{array}\pmod{1129}.
 \tag{E.70.12}
\]
At $t=828$ the radicand in (E.70.1) is $1115$ and has
exactly the eight roots
\[
 35,44,235,511,618,894,1085,1094.
\]
Formula (E.70.4)
gives $x=28,941,1101,188$, each twice. Their two norm
residues are $(528,1040)$ or $(89,601)$, all nonsquares.
Consequently every local sheet has squareclass vector $(1,1)$,
whereas every global class in $\mathcal S$ has vector $(0,0)$.
All these entries and all root derivatives are units. Simple-root
Hensel lifting therefore preserves both the eight sheets and
their nonzero squareclasses throughout the stated residue ball.
The program \path{code/check_genus_two_norm.py} independently
checks these finite assertions, all sixteen supported candidates,
the birational identities, both geometric groups, and their
quotient map; its complete output is retained.

Finally, a section in combined class $[-24]$ trivializes the
pullback of (E.70.1), and choosing one fixed sheet restricts it
to the arithmetic cover subgroup. Map it through (E.70.5) and
(E.70.7). This yields one continuous global class in
$H^1(G_\Q,\overline M)$. If the section is locally point-or-cusp
matching, each local geometric conjugator transports that same
sheet to a rational local cover point, or to a rational boundary
point after compactification. Its Jacobian class belongs to the
local Kummer image, at the real place as well. Thus the global
class lies in $\mathcal S$. Every possible transported sheet
at $1129$ was excluded above. This proves the asserted empty
locally matching locus, with no choice of sheet suppressed.
\end{proof}

\begin{result}{E.71}{Proved}{An original factorial ghost first fails local matching at class two}
There is an output of the literal recursion E.102,
\[
 A=3\prod_j q_j^{(j+4)!},\qquad
 B=-2\prod_j r_j^{(j+4)!},
 \tag{E.71.1}
\]
whose ordinary localization at $1129$ has centre in
$828+1129\Z_{1129}$. Its abelian section is locally a point
section everywhere, including the real place. Its everywhere
locally point-or-cusp-matching integral class-two fibre is empty,
while its unrestricted global class-two fibre has the twenty-four
classes of A.27. Thus its first local-matching failure is exactly
class two. This confirms GT\_TOWER\_CHECK C.2.2 in its stated
local-matching scope and does not assign this pair a first
unrestricted nilpotent failure.
\end{result}
\begin{proof}
In the first old-prime local choice of E.102 at $1129$, use
$X=2$ and the root $Y\equiv56\pmod{1129}$ of
$3X^{24}-2Y^{24}=1$. Indeed
\[
 2^{24}=276,\quad56^{24}=978,\quad
 3\cdot276-2\cdot978=1\pmod{1129}.
 \tag{E.71.2}
\]
The derivative in $Y$ is a unit, so this residue root has a unique
Hensel lift. This prime is among the first processed primes and
does not divide the next exponent $120$; hence the choice satisfies
every local requirement of that construction. Make the original
CRT approximations to the exact root at their specified precision.
The first updated coefficient has residue $3\cdot276=828$.
At every later stage take $X=1$ and the $Y$-root near one.
The maintained error bound in E.102 exceeds the Hensel threshold
and gives the required depth even after $1129$ divides the
factorial. Every later multiplier at $1129$ is consequently one
modulo $1129$. All other choices, their splitting conditions,
and the all-place local limit proof are unchanged.

All update exponents are multiples of eight. Hence
$(A,B)\bmod8=([3],[-2])$ and $AB^3\bmod8=[-24]$.
The supports remain disjoint and infinite. Fixing the least
admissible CRT and prime choices elsewhere specifies an effective
recursion; we do not claim to have searched its enormous first
prime. The finite identities actually used in its first local
choice were executed by the retained program.

A locally matching integral class-two section would map to the
empty locus of E.70. Its local ordinary centre must be the stated
one by full local Kummer injectivity. A cusp cannot replace it:
the three abelian cusp patterns would require respectively
$1-t=1$, $t=1$, or $(1-t)/t=-1$ in the local field, none of
which holds. E.70 therefore excludes it. Ordinariness and the
real sign pair $(+,-)$ give an everywhere matching abelian section,
so there was no earlier failure. The unrestricted cup vanishes by
A.25, and A.27 gives its unrestricted fibre; these are compatible
with the proved local-matching exclusion.
\end{proof}

\begin{result}{E.72}{Proved}{The separating prime gives no local third obstruction}
For every unrestricted global class-two section above the pair of
E.71, its local class-three obstruction at $q=1129$ is zero for
every coefficient prime. For each $\ell\ne1129$, that specified
local class-two section lifts to a continuous section of the
entire geometrically pro-$\ell$ local extension over $\Q_{1129}$.
This confirms GT\_TOWER\_CHECK C.2.4. In particular the finite
global Selmer exclusion in E.70 is not an unrestricted local
third obstruction.
\end{result}
\begin{proof}
Both local abelian coordinates are units. For $\ell\ne q$ their
inertia values vanish. Any remaining class-two inertia is central
in the torsion-free module $\Z_\ell(2)$. The relation
$FTF^{-1}=T^q$ and Frobenius weight $q^2$ force
$(q^2-q)h(T)=0$, hence $h(T)=0$. Wild and prime-to-$\ell$
tame inertia also map trivially. Thus the specified section is
unramified. E.32 proves that the actual tangential action is
unramified in every nilpotent class at $q\ne\ell$; the inverse
limit is the full free pro-$\ell$ group, since every finite
$\ell$-group has finite nilpotency class.

The unramified arithmetic extension is therefore
$\Pi^{(\ell)}\rtimes\widehat\Z$. Choose any lift of the specified
Frobenius element from class two to this profinite group. The
universal property of $\widehat\Z$ extends its powers to a
continuous homomorphism. Its projection sends the generator to
the generator and is the identity. Inflation gives the desired
local section with exactly the prescribed class-two reduction.

For $\ell=q=1129$, finite local Tate duality
\cite[Corollary~I.2.3, p.~28]{CohMilneADT} identifies the dual
of $H^2(G_{\Q_q},\Z/q^n(3))$ with
$H^0(G_{\Q_q},\Z/q^n(-2))$. The mod-$q$ cyclotomic image
has order $q-1$, which does not divide two. Thus one scalar has
$\chi^{-2}-1$ a unit; all these invariant groups vanish.
The local finite $H^1$ groups are finite, so their inverse-system
derived limit vanishes. Continuous $H^2(G_{\Q_q},\Z_q(3))$
is consequently zero as well. This treats the coefficient prime
without asserting any local point-matching property for the lift.
\end{proof}
